\chapter{Literature Review}
\label{ch:review}

This chapter summarises the key evidence that shaped the design of this study.
A full critical review of fifty sources is provided in the companion document
\textit{Statistical and Machine-Learning Approaches to Real-Time Sepsis Onset
Prediction on MIMIC-IV: An Extended Critical Literature Review}. This chapter
extracts the seven evidence strands that directly informed the study's design
choices.

\subsection*{Key Concepts}

Several domain-specific terms recur throughout this review:

\begin{itemize}
  \item \textbf{MIMIC-IV} (Medical Information Mart for Intensive Care, version
    IV): a freely available, de-identified database of approximately 74,000 ICU
    admissions at Beth Israel Deaconess Medical Center, Boston, covering
    2008--2019 \cite{johnson2023mimiciv}.
  \item \textbf{Sepsis-3}: the current international consensus definition of
    sepsis as life-threatening organ dysfunction caused by a dysregulated host
    response to infection, operationalised through a combination of suspected
    infection (antibiotic order plus culture draw) and a rise in the Sequential
    Organ Failure Assessment (SOFA) score of $\geq$2 points
    \cite{singer2016sepsis3}.
  \item \textbf{AUROC} (Area Under the Receiver Operating Characteristic
    curve): a discrimination metric that measures a model's ability to rank
    positive cases above negative cases, ranging from 0.5 (random) to 1.0
    (perfect separation). It does not account for alert timing or
    false-alarm cost.
  \item \textbf{Utility Score}: a net-benefit metric proposed by the
    PhysioNet/Computing in Cardiology Challenge that rewards early, true
    alerts and penalises late alerts and false alarms, producing values that
    can be negative when a model causes more harm than benefit.
  \item \textbf{NEWS2} (National Early Warning Score 2): a standardised
    clinical scoring system aggregating six physiological parameters (respiratory
    rate, oxygen saturation, systolic blood pressure, pulse rate, level of
    consciousness, and temperature) into a single risk score used for bedside
    deterioration screening.
\end{itemize}

\subsection*{Research Questions}

This literature review is structured around four questions that directly
inform the study's design:

\begin{enumerate}
  \item[\textbf{RQ1.}] How much does the choice of sepsis label definition
    affect model performance, and is this effect comparable in magnitude to
    the choice of model architecture?
  \item[\textbf{RQ2.}] What methodological pitfalls---validation strategy,
    metric selection, treatment-time leakage---most distort reported
    performance in the sepsis prediction literature?
  \item[\textbf{RQ3.}] What statistical framework is appropriate for
    modelling hourly ICU event data when the standard proportional-hazards
    assumption is violated?
  \item[\textbf{RQ4.}] Does the existing equity literature on sepsis
    prediction account for label-level variation across demographic subgroups,
    or only model-level variation?
\end{enumerate}

\subsection*{Current State of the Field}

Real-time sepsis prediction has attracted substantial research effort over
the past decade. Wang et al.\ \cite{wang2025srpm} identified ninety-one
models published between 2015 and 2024, spanning logistic regression,
gradient-boosted trees, recurrent neural networks, and transformer
architectures. Despite this volume, most studies share three weaknesses:
they use a single label definition, report AUROC as the primary metric,
and validate on a random patient-level split. The evidence reviewed in
this chapter shows that each of these choices individually inflates or
obscures model performance, and their combination produces a literature
in which high reported AUROCs coexist with models that fail to improve
clinical outcomes at external sites. Table~\ref{tab:study_comparison}
provides a structured comparison of the key studies that inform this
analysis.

\begin{table}[H]
  \centering\small
  \begin{tabularx}{\textwidth}{p{2.4cm}p{1.6cm}p{2.2cm}X p{2.2cm}}
    \toprule
    \textbf{Study} & \textbf{Dataset} & \textbf{Method} &
    \textbf{Key Finding} & \textbf{Gap Addressed} \\
    \midrule
    Johnson et al.\ \cite{johnson2018comparative} &
    MIMIC-III &
    5 label algorithms &
    Agreement between sepsis definitions only moderate ($\alpha$\,=\,0.40--0.62) &
    Label ambiguity \\
    \addlinespace
    Cohen et al.\ \cite{cohen2024subtle} &
    MIMIC-III &
    Trees, DL, survival &
    Label change $\approx$ model change in performance impact &
    Label as variable \\
    \addlinespace
    Dutta et al.\ \cite{dutta2026performance} &
    9 hospitals (198k) &
    Commercial GBT &
    Label sensitivity holds prospectively at scale &
    External validity \\
    \addlinespace
    Kamran et al.\ \cite{kamran2024evaluation} &
    MIMIC-IV &
    Multiple models &
    AUROC inflated when treatment-time leakage not excluded &
    Treatment bias \\
    \addlinespace
    Guo et al.\ \cite{guo2022temporal} &
    MIMIC-IV &
    Domain adaptation &
    Sepsis worst-affected by temporal shift; random split inflates &
    Validation strategy \\
    \addlinespace
    Wang et al.\ \cite{wang2025srpm} &
    91-study review &
    Meta-analysis &
    AUROC survives external shift; Utility Score collapses &
    Metric selection \\
    \addlinespace
    Wang et al.\ \cite{wang2022equity} &
    MIMIC-III &
    Subgroup analysis &
    Performance varies by race, ethnicity, insurance, language &
    Label-by-subgroup gap \\
    \bottomrule
  \end{tabularx}
  \caption{Comparison of key studies reviewed, their datasets, methods, and the design gaps each addresses.}
  \label{tab:study_comparison}
\end{table}

\section{The Sepsis Label Is Not a Single Object}

The common assumption that ``the MIMIC-IV sepsis cohort'' is one well-defined
group is wrong. Two studies show why.

Johnson et al.\ \cite{johnson2018comparative} applied five different sepsis
identification algorithms to 11,791 ICU admissions in MIMIC and found only
moderate agreement between them (Cronbach's $\alpha$ = 0.40--0.62). The
implication is clear: ``the MIMIC sepsis cohort'' is not one group but a family
of related but different groups, each created by a specific reading of Sepsis-3.

Cohen et al.\ \cite{cohen2024subtle} showed what this means for prediction.
On MIMIC-III, using three different but reasonable Sepsis-3 interpretations, the
identified cohort ranged from 867 to 2,178 admissions. Across multiple model
types (tree-based, deep learning, and survival models), changing the label
definition moved performance by as much as changing the model itself. This is
the most important finding behind the present study: it means the label must be
treated as a primary variable, not a fixed background choice.

Dutta et al.\ \cite{dutta2026performance} confirmed this is not a
MIMIC-specific artefact. Using a commercial gradient-boosted model across nine
hospitals and 198,494 encounters, they found that performance varied
substantially depending on whether Sepsis-3, SEP-1, or the CDC Adult Sepsis
Event definition was used. The label sensitivity finding holds prospectively,
at scale, with a real deployed model.

Figure~\ref{fig:label_variance_concept} illustrates this core problem: the
same patient population, subjected to different but equally valid Sepsis-3
interpretations, produces overlapping but distinct cohorts, each yielding
different model performance.

\begin{figure}[H]
  \centering
  \begin{tikzpicture}[
    every node/.style={font=\small},
    cohort/.style={draw, circle, minimum size=2.8cm, thick, fill opacity=0.15},
  ]
    \node[cohort, fill=ATUGreen, draw=ATUGreen] (a) at (0, 0) {};
    \node[cohort, fill=ATUOrange, draw=ATUOrange] (b) at (1.6, 0) {};
    \node[cohort, fill=ATUNavy, draw=ATUNavy] (c) at (0.8, 1.3) {};

    \node[below=0.2cm, font=\footnotesize\bfseries, text=ATUGreen] at (a.south) {Label A};
    \node[below=0.2cm, font=\footnotesize\bfseries, text=ATUOrange] at (b.south) {Label B};
    \node[above=0.2cm, font=\footnotesize\bfseries, text=ATUNavy] at (c.north) {Label C};

    \node[font=\scriptsize] at (0.8, 0.4) {\textbf{Shared}};

    \draw[decorate, decoration={brace, amplitude=5pt, raise=3pt}]
      (-1.6, -2.0) -- (3.2, -2.0)
      node[midway, below=10pt, font=\footnotesize, align=center]
      {Same ICU population, same EHR data\\different ``sepsis cohorts''};

    \node[draw, rounded corners, fill=white, inner sep=6pt,
          font=\footnotesize, align=left] at (6.5, 0.5) {
      \textbf{Impact on AUROC:}\\[2pt]
      Label A: 0.82\\
      Label B: 0.78\\
      Label C: 0.85\\[4pt]
      $\Delta_{\text{label}} \approx \Delta_{\text{model}}$
    };
  \end{tikzpicture}
  \caption{Conceptual illustration of label variance: three valid Sepsis-3
    interpretations applied to the same ICU population produce overlapping but
    distinct cohorts, each yielding different model performance. The
    performance gap between labels is comparable to the gap between model
    architectures.}
  \label{fig:label_variance_concept}
\end{figure}

\section{The Label Depends on Clinical Action}

The Sepsis-3 label is built from antibiotic orders, culture draws, and SOFA
score rises. Kamran et al.\ \cite{kamran2024evaluation} observed that
clinicians often recognise and begin treating sepsis before the formal criteria
are met. A model trained to predict the criteria-meeting time therefore partly
learns to predict when the clinician will act. They evaluated performance in
the window before treatment started and showed that apparent AUROC was
substantially inflated without that restriction. The model had learned clinician
suspicion, not independent physiological deterioration.

Lauritsen et al.\ \cite{lauritsen2021framing} formalised the related
concept of \textit{framing}: the choice of observation window, prediction
horizon, and event trigger. They showed that framing changes can reverse the
sign of a feature's effect; their example is SpO$_2$. This means a hazard
ratio whose sign depends on how the study was framed is not a clinically
meaningful quantity.

\section{The Validation Split Must Be Time-Based}

Guo et al.\ \cite{guo2022temporal} tested four clinical prediction tasks
on MIMIC-IV under temporal shift (training on older data, testing on newer data).
Sepsis was the worst-affected task. Domain-generalisation and domain-adaptation
techniques did not help. A random patient-level split therefore inflates
performance more for sepsis than for any other clinical outcome.

This finding drives a key design decision: the primary validation in this study
uses a time-based split (train on 2008--2016, test on 2017--2019). A random
split is only used to measure how much performance it inflates.

\section{AUROC Hides the Clinical Failure}

Wang et al.\ \cite{wang2025srpm} reviewed ninety-one real-time sepsis
models. Median internal AUROC was 0.811, external AUROC was 0.783, a small,
insignificant drop. But median internal Utility Score was $+0.381$, external
was $-0.164$: the models went from clinically useful to actively harmful. AUROC
survived this transition because it only measures whether the model can rank
patients by risk. The Utility Score failed because it measures whether alerts
arrive in time and at acceptable false-alert rates, and at external sites,
they did not. Reporting AUROC as the main metric does not just miss the failure;
it actively hides it.

\section{The Survival Model for Hourly ICU Data}

The standard survival approach uses a Cox proportional-hazards model. On
MIMIC-IV, the proportional-hazards assumption is rejected (Schoenfeld test,
$p < 10^{-16}$). D'Agostino et al.\ \cite{dagostino1990pooled} provided
the solution: pool observations across short time intervals and use logistic
regression for each interval. This pooled logistic regression is mathematically
equivalent to time-dependent Cox regression when intervals are short and
per-interval event rates are low. One-hour ICU intervals with sepsis onset
probability below 1\% meet both conditions.

Resche-Rigon et al.\ \cite{rescherigon2006competing} and Heyard et al.\
\cite{heyard2019dynamic, heyard2020evaluation} added the competing-risks
extension. In the ICU, patients can leave alive or die without developing
sepsis. Treating these events as simple censoring overstates sepsis incidence
and biases model estimates. The multinomial cause-specific formulation used
in this study handles this correctly.

\section{NEWS2 as the Baseline Comparator}

Most MIMIC-IV sepsis models compare against SIRS, qSOFA, and NEWS. Evans et al.\
\cite{evans2021ssc} issued a strong recommendation against using qSOFA
alone for screening due to poor sensitivity. Mellhammar et al.\
\cite{mellhammar2020news2} showed that a LASSO-derived score could not
beat NEWS2. The current clinical guidelines have therefore already retired qSOFA
as a standalone screening tool, making NEWS2 the minimum meaningful comparator.
This study retains SIRS and qSOFA for completeness but uses NEWS2 as the primary
rule-based benchmark.

\section{Equity Analysis: Performance and Label Sensitivity Together}

Wang, Li, Naidech and Luo \cite{wang2022equity} showed that sepsis model
performance varies significantly across racial, ethnic, language, and insurance
subgroups. However, they did not examine whether the label itself assigns sepsis
status differently across those groups. This study adds that missing piece: for
each subgroup, it reports both the variation in who gets labelled as septic
under each label variant and the variation in model performance. This
combination has not been reported for any sepsis model.

\section{Summary: Design Decisions and Their Evidence}

Table~\ref{tab:design_warrant} maps each major design decision to the evidence
that supports it.

\begin{table}[H]
  \centering\small
  \begin{tabular}{p{3.2cm}p{6cm}p{3.5cm}}
    \toprule
    \textbf{Decision} & \textbf{Reason} & \textbf{Source} \\
    \midrule
    Three label variants & Label variance $\approx$ model variance &
      Cohen et al.\ \cite{cohen2024subtle} \\
    Label as variable & Variance is a finding, not a nuisance &
      Cohen; Dutta et al.\ \cite{dutta2026performance} \\
    Treatment-anchored eval & Label depends on clinical action &
      Kamran et al.\ \cite{kamran2024evaluation} \\
    Time-based split & Sepsis worst-affected by temporal shift &
      Guo et al.\ \cite{guo2022temporal} \\
    Utility Score primary & AUROC hides clinical failure &
      Wang et al.\ \cite{wang2025srpm} \\
    Calibration co-primary & Critical for clinical usefulness &
      Van Calster et al.\ \cite{vancalster2019calibration} \\
    Discrete-time model & Cox-equivalent; no PH assumption &
      D'Agostino \cite{dagostino1990pooled};
      Heyard \cite{heyard2019dynamic} \\
    NEWS2 comparator & qSOFA no longer recommended &
      Evans et al.\ \cite{evans2021ssc} \\
    GBT comparator & Cross-study comparison needs same label &
      Cohen et al.\ \cite{cohen2024subtle} \\
    Subgroup label sens. & Equity literature omits label-by-subgroup &
      Wang et al.\ \cite{wang2022equity} \\
    \bottomrule
  \end{tabular}
  \caption{Design decisions and their evidence warrants.}
  \label{tab:design_warrant}
\end{table}
